\begin{ouverture}[vecteursplan]{15}{Translations et vecteurs}{Partie 4 · Espace et géométrie}
{Symétries · Translations · Vecteurs égaux · Somme · Relation de Chasles}
\objectif{Construire l’image d’une figure par symétrie, demi-tour ou translation.}
\objectif{Reconnaître des vecteurs égaux et un parallélogramme.}
\objectif{Construire une somme de vecteurs et utiliser la relation de Chasles.}
\prerequis{Symétrie axiale, demi-tour, translation et parallélogrammes.}
\end{ouverture}
\pagestyle{otMain}\essentiel\fontsize{11}{15}\selectfont
\essSection{1}{Translations et vecteurs}
\begin{definition}La translation qui transforme A en B transforme tout point C en l’unique point D tel que ABDC soit un parallélogramme, éventuellement aplati. On l’appelle la translation de vecteur $\vect{AB}$.\end{definition}
\begin{definition}Pour deux points distincts A et B, le vecteur $\vect{AB}$ est défini par la direction de la droite (AB), le sens de A vers B et la longueur AB, appelée norme. On le représente par une flèche ; deux vecteurs sont égaux s’ils ont même direction, même sens et même longueur.\end{definition}
\begin{propriete}$\vect{AB}=\vect{CD}$ si et seulement si ABDC est un parallélogramme, éventuellement aplati. Deux vecteurs égaux définissent la même translation.\end{propriete}
\begin{definition}Le vecteur nul est $\vect{AA}=\vect{BB}=\vec0$ : la translation correspondante ne déplace aucun point. L’opposé de $\vect{AB}$ est $-\vect{AB}=\vect{BA}$, de même direction et même longueur, mais de sens contraire.\end{definition}
\begin{attention}L’ordre des lettres compte : $\vect{AB}=\vect{DC}$ signifie que ABCD est un parallélogramme, alors que $\vect{AB}=\vect{CD}$ correspond à ABDC. Symétries, demi-tours et translations conservent longueurs, angles, alignements et aires.\end{attention}
\clearpage
\essSection{2}{Somme de vecteurs et relation de Chasles}
\begin{propriete}Enchaîner la translation de vecteur $\vect{AB}$ puis celle de vecteur $\vect{BC}$ revient à appliquer la translation de vecteur $\vect{AC}$. On définit ainsi la somme :
\[\vect{AB}+\vect{BC}=\vect{AC}\qquad\text{(relation de Chasles).}\]
En particulier $\vect{AB}+\vect{BA}=\vect{AA}=\vec0$.\end{propriete}
\begin{savoirfaire}[Construire $\vec u+\vec v$]À partir d’un point O, tracer un représentant $\vect{OP}$ de $\vec u$, puis, à partir de P, un représentant $\vect{PS}$ de $\vec v$. Alors $\vec u+\vec v=\vect{OS}$. On note $2\vect{AB}$ la somme $\vect{AB}+\vect{AB}$ : même direction, même sens, longueur double.\end{savoirfaire}
\begin{propriete}[Règle du parallélogramme]Si ABCD est un parallélogramme, alors $\vect{AB}+\vect{AD}=\vect{AC}$ : on remplace $\vect{AD}$ par le vecteur égal $\vect{BC}$, puis on applique Chasles.\end{propriete}
\begin{attention}La somme de deux vecteurs n’est pas la somme de leurs longueurs : $AB+BC$ n’est égal à $AC$ que si B appartient au segment [AC].\end{attention}
\clearpage\rubrique{Découverte}
\begin{blocLivre}{ex-15-1}\begin{exercice}[1]{Symétrie et demi-tour}
Sur le quadrillage, construire l’image A’B’C’ du triangle ABC par la symétrie d’axe (d), puis son image A’’B’’C’’ par le demi-tour de centre O. Ces transformations conservent-elles les longueurs ?
\label{fig-ex-15-1}\figureGeo{vect-01}
\end{exercice}\end{blocLivre}
\begin{blocLivre}{ex-15-2}\begin{exercice}[1]{Caractériser un vecteur}
Le côté d’un carreau représente 1 cm. Décrire le vecteur $\vect{AB}$ : direction, sens et longueur. Construire ensuite le point D tel que $\vect{CD}=\vect{AB}$.
\label{fig-ex-15-2}\figureGeo{vect-02}
\end{exercice}\end{blocLivre}
\begin{blocLivre}{ex-15-3}\begin{exercice}[1]{Des vecteurs égaux}
Parmi les vecteurs $\vect{CD}$, $\vect{EF}$, $\vect{GH}$ et $\vect{IJ}$ dessinés, lesquels sont égaux à $\vect{AB}$ ? Lequel est son opposé ?
\label{fig-ex-15-3}\figureGeo{vect-03}
\end{exercice}\end{blocLivre}
\begin{blocLivre}{ex-15-4}\begin{exercice}[1]{Image par une translation}
Construire l’image D du point C par la translation de vecteur $\vect{AB}$. Quelle est la nature du quadrilatère ABDC ?
\label{fig-ex-15-4}\figureGeo{vect-04}
\end{exercice}\end{blocLivre}
\begin{blocLivre}{ex-15-5}\begin{exercice}[1]{Vecteur nul et opposé}
Simplifier $\vect{AA}$. Donner un autre nom du vecteur $-\vect{MN}$, puis calculer $\vect{MN}+\vect{NM}$.
\label{fig-ex-15-5}\figureGeo{vect-05}
\end{exercice}\end{blocLivre}

\clearpage\rubrique{Entraînement}
\begin{blocLivre}{ex-15-6}\begin{exercice}[2]{Parallélogramme et égalités}
ABCD est un parallélogramme. Compléter avec des points de la figure : $\vect{AB}=\vect{D\,\ldots}$ ; $\vect{AD}=\vect{\ldots\,C}$ ; $\vect{BA}=\vect{C\,\ldots}$ ; $\vect{CB}=\vect{\ldots\,A}$.
\label{fig-ex-15-6}\figureGeo{vect-06}
\end{exercice}\end{blocLivre}
\begin{blocLivre}{ex-15-7}\begin{exercice}[2]{Démontrer un parallélogramme}
On sait que $\vect{EF}=\vect{HG}$. Démontrer que EFGH est un parallélogramme, puis vérifier cette égalité sur le quadrillage.
\label{fig-ex-15-7}\figureGeo{vect-07}
\end{exercice}\end{blocLivre}
\begin{blocLivre}{ex-15-8}\begin{exercice}[2]{La relation de Chasles}
Simplifier : $\vect{AB}+\vect{BC}$ ; $\vect{MN}+\vect{NP}+\vect{PQ}$ ; $\vect{AB}+\vect{CA}$ ; $\vect{AB}+\vect{BA}$.
\label{fig-ex-15-8}\figureGeo{vect-08}
\end{exercice}\end{blocLivre}
\begin{blocLivre}{ex-15-9}\begin{exercice}[2]{Construire une somme}
Les vecteurs $\vec u$ et $\vec v$ sont dessinés ; le côté d’un carreau représente 1 cm. Construire un représentant de $\vec u+\vec v$ d’origine O, puis calculer sa longueur.
\label{fig-ex-15-9}\figureGeo{vect-09}
\end{exercice}\end{blocLivre}
\begin{blocLivre}{ex-15-10}\begin{exercice}[2]{Deux translations successives}
On applique au point M la translation de vecteur $\vect{AB}$, ce qui donne M’, puis la translation de vecteur $\vect{BC}$, ce qui donne M’’. Construire M’ et M’’. Par quelle translation unique passe-t-on de M à M’’ ?
\label{fig-ex-15-10}\figureGeo{vect-10}
\end{exercice}\end{blocLivre}

\clearpage\rubrique{Problèmes}
\begin{blocLivre}{ex-15-11}\begin{exercice}[3]{La règle du parallélogramme}
ABCD est un parallélogramme. Démontrer que $\vect{AB}+\vect{AD}=\vect{AC}$.
\label{fig-ex-15-11}\figureGeo{vect-11}
\end{exercice}\end{blocLivre}
\begin{blocLivre}{ex-15-12}\begin{exercice}[3]{Milieu et vecteurs}
I est le milieu de [AB]. Démontrer que $\vect{AI}=\vect{IB}$ et que $\vect{AB}=2\vect{AI}$. Réciproquement, si un point J vérifie $\vect{AJ}=\vect{JB}$, que peut-on dire de J ?
\label{fig-ex-15-12}\figureGeo{vect-12}
\end{exercice}\end{blocLivre}
\begin{blocLivre}{ex-15-13}\begin{exercice}[3]{Un drone}
Un drone part de D, se déplace de 60 m vers l’est, puis de 80 m vers le nord. Représenter les deux déplacements par des vecteurs sur le quadrillage, où un carreau représente 20 m. Quel déplacement unique produit le même effet ? Quelle distance sépare le point d’arrivée F du point de départ ?
\label{fig-ex-15-13}\figureGeo{vect-13}
\end{exercice}\end{blocLivre}
\begin{blocLivre}{ex-15-14}\begin{exercice}[3]{Une frise}
Une frise est obtenue en répétant un motif par la translation de vecteur $\vec u$, de longueur 4 cm. Quelle translation transforme le motif 1 en motif 6 ? Le motif 6 en motif 2 ? Donner les longueurs des vecteurs correspondants.
\label{fig-ex-15-14}\figureGeo{vect-frise}
\end{exercice}\end{blocLivre}
\begin{blocLivre}{ex-15-15}\begin{exercice}[3]{Démontrer avec des vecteurs}
ABCD est un parallélogramme et E est le point tel que $\vect{BE}=\vect{AB}$. Démontrer que DCEB est un parallélogramme.
\label{fig-ex-15-15}\figureGeo{vect-15}
\end{exercice}\end{blocLivre}

\clearpage
